\section{Semantics}
\label{sec:semantics}

We now build the structures against which soundness is proved: a
proof-irrelevant, classical, syntactic semantics of the source, and a
first-order structure $\ML$ over the same individuals interpreting
$\SigFOL$.  Throughout, $\Env$ is a fixed well-formed definitional
environment with $\Sfrag$ declarations.  The construction follows
\cite{Extraction2026}; the new material is the interpretation of indexed
families as sets of (index tuple, element) pairs, the \emph{index
reflection} lemma --- the semantic content of fording --- and element
ranks.

\subsection{The domain}
\label{sec:semantics-domain}

\begin{definition}[Domain]
\label{def:domain}
Let $\Dom$ be the set of closed $\lbox$ terms modulo $\lconv$, written
$\cls{e}$ for classes.  $\Dom$ is non-empty ($\cls{\bx} \in \Dom$) and
carries the application $d \cdot d' \defeq \cls{e\,e'}$ for
representatives, well defined since $\lconv$ is a congruence.
\end{definition}

$\Dom$ deliberately contains \emph{junk}: classes such as
$\cls{\csO\,\csO}$ or $\cls{\mathsf{fix}\,f.f}$ that are the value of no
source term of any type, and --- new in the indexed setting ---
\emph{well-formed constructor forms at the wrong index}, such as
$\cls{\csvnil}$ considered at index $\cls{\csS(\csO)}$.  By
\cref{cor:discrimination}, distinct constructor forms lie in distinct
classes and junk like $\cls{\csO\,\csO}$ lies outside every constructor
form.  Junk is what makes the guard analysis of \cref{sec:formulas}
non-trivial; keeping it in the model is what lets the model arbitrate
which axioms need guards, and the wrong-index phenomenon is what makes
the \emph{index} arguments of guards non-negotiable
(\cref{prop:index-dropping}).

\begin{definition}[Valuations and term denotation]
\label{def:valuation}
A \emph{valuation} $\rho$ assigns to finitely many term variables
elements of $\Dom$ and to finitely many predicate variables $X$ (of kind
$(\sigma_1, \dots, \sigma_n) \to \Prop$) functions $\Dom^n \to \{0,1\}$.
Proof variables are not assigned anything.  For a $\CICi$ term $t$ whose
free term variables are assigned by $\rho$,
\[
  \sem{t}_\rho \defeq \cls{\erase{t}^\rho} \in \Dom,
\]
where $\erase{t}^\rho$ substitutes a representative of $\rho(x)$ for
each free variable $x$ of $\erase{t}$; this is well defined exactly as
in \cite{Extraction2026} (free variables of $\erase{t}$ are free term
variables of $t$; classes are representative-independent).  For an index
tuple we write $\sem{\tup{u}}_\rho = (\sem{u_1}_\rho, \dots,
\sem{u_m}_\rho) \in \Dom^m$.
\end{definition}

\subsection{Interpretation of types, propositions and predicates}
\label{sec:semantics-interp}

\begin{definition}[Interpretations]
\label{def:interp}
By recursion on the position of declarations in $\Env$ and, within each
declaration, on the structure of phrases, we define for every set type
$\sigma$, proposition $\varphi$ and predicate expression $P$ over a
prefix of $\Env$, and every valuation $\rho$ covering their free
variables:
a set $\semr{\sigma}{\rho} \subseteq \Dom$; a truth value
$\pvr{\varphi}{\rho} \in \{0,1\}$; a function $\kv{P}_\rho : \Dom^n \to
\{0,1\}$.

\emph{Indexed datatypes.}  For $I : (\sigma_1, \dots, \sigma_m) \to
\Set$ with constructors $\wcns{c}_i : \forall \tup{y}{:}\tup{A}_i.\
\varphi_{i1} \Rightarrow \dots \Rightarrow \varphi_{i\ell_i} \Rightarrow
I(\tup{t}_i)$ ($r_i$ informative arguments), define stages $\stage{I}{s}
\subseteq \Dom^m \times \Dom$ by $\stage{I}{0} = \emptyset$ and
\begin{multline*}
  \stage{I}{s+1} = \Bigl\{\, \bigl(\sem{\tup{t}_i}_{\rho_0},\,
    \cls{\wcns{c}_i(e_1, \dots, e_{r_i})}\bigr) \;\Bigm|\;
    i \le M,\ \rho_0 = [\tup{y}_i \mapsto \cls{\tup{e}}\,],\\
    \cls{e_j} \in \semr{A_{ij}}{\rho_0} \text{ for non-recursive } j,\\
    \bigl(\sem{\tup{v}_{ij}}_{\rho_0}, \cls{e_j}\bigr) \in \stage{I}{s}
    \text{ for recursive } j \text{ with } A_{ij} = I(\tup{v}_{ij}),\\
    \pvr{\varphi_{il}}{\rho_0} = 1 \text{ for all } l \,\Bigr\},
\end{multline*}
which is well defined: non-recursive argument types and all premises
are $I$-free phrases over the environment prefix and the telescope, so
their interpretations and truth values are supplied by the outer
recursion, and index-term denotations $\sem{\cdot}_{\rho_0}$ are erasure
classes, defined independently of any type interpretation.  The stages
are increasing in $s$ (induction on $s$, the defining condition being
monotone in the stage at recursive positions).  Set $\semr{I}{} =
\bigcup_s \stage{I}{s}$,\footnote{A reader may expect the invariant
$\semr{I}{} \subseteq \semr{\sigma_1}{} \times \cdots \times
\semr{\sigma_m}{} \times \Dom$ --- index tuples inhabiting their index
types.  It is not built into the stages, and it is stated and needed
nowhere in the development: index-term denotations are erasure
classes, supplied independently of every type interpretation, and no
lemma ever requires an index tuple to inhabit anything.  This is
itself information about the translation: no guard is ever emitted for
an index argument place (\cref{def:theory}), and the model shows why
none is missed.} and for an occurrence:
\[
  \semr{I(\tup{u})}{\rho} =
  \{\, d \in \Dom \mid (\sem{\tup{u}}_\rho,\, d) \in \semr{I}{} \,\}.
\]

\emph{Other set types} (as in \cite{Extraction2026}):
\begin{align*}
  \semr{\Pi x{:}\sigma.\tau}{\rho} &=
    \{\, d \in \Dom \mid \forall e \in \semr{\sigma}{\rho}.\;
      d \cdot e \in \semr{\tau}{\rho[x \mapsto e]} \,\}\\
  \semr{\varphi \to \tau}{\rho} &=
    \{\, d \in \Dom \mid \pvr{\varphi}{\rho} = 1 \implies d \in
      \semr{\tau}{\rho} \,\}\\
  \semr{\refty{x{:}\sigma}{\varphi}}{\rho} &=
    \{\, d \in \semr{\sigma}{\rho} \mid \pvr{\varphi}{\rho[x \mapsto d]}
      = 1 \,\}
\end{align*}

\emph{Inductive predicates.}  For $J : (\sigma_1, \dots, \sigma_k) \to
\Prop$ with clauses $\wcns{k}_j : \forall \tup{x}{:}\tup{\tau}_j.\
\varphi_{j1} \to \dots \to \varphi_{j\ell_j} \to J(\tup{u}_j)$, stages
$\pv{J}^0 = \emptyset \subseteq \Dom^k$ and
\begin{multline*}
  \pv{J}^{s+1} = \Bigl\{\, \sem{\tup{u}_j}_{\rho_0} \;\Bigm|\;
    j,\ \rho_0 = [\tup{x} \mapsto \tup{d}\,],\
    d_i \in \semr{\tau_{ji}}{\rho_0},\\
    \pvr{\varphi}{\rho_0} = 1 \text{ for non-recursive premises},\\
    \sem{\tup{v}}_{\rho_0} \in \pv{J}^s \text{ for recursive premises }
    J(\tup{v}) \,\Bigr\},
\end{multline*}
with $\pv{J} = \bigcup_s \pv{J}^s$.

\emph{Propositions and predicate expressions}: exactly the clauses of
\cite{Extraction2026}: classical truth tables for the connectives,
$\min/\max$ over $\semr{\sigma}{\rho}$ for quantifiers (with $\min
\emptyset = 1$, $\max \emptyset = 0$), $\min$ over all $Q : \Dom^n \to
\{0,1\}$ for predicate quantification, $\pvr{a =_\sigma b}{\rho} = 1$
iff $\sem{a}_\rho = \sem{b}_\rho$, $\pvr{J(\tup{a})}{\rho} = 1$ iff
$\sem{\tup{a}}_\rho \in \pv{J}$, $\pvr{X(\tup{a})}{\rho} =
\rho(X)(\sem{\tup{a}}_\rho)$, $\kv{X}_\rho = \rho(X)$ and
$\kv{\lambda(\tup{x}{:}\tup{\sigma}).\varphi}_\rho(\tup{d}) =
\pvr{\varphi}{\rho[\tup{x} \mapsto \tup{d}\,]}$.
\end{definition}

The recursion is well founded exactly as before: stage constructions
refer only to earlier declarations and structurally smaller phrases of
their own declaration; all other clauses recurse on proper sub-phrases.
The semantics is classical and proof-irrelevant by construction: truth
values do not depend on proofs, and $\sem{\cdot}$ ignores proof
subterms.  The interpretation of equality is untyped (class equality,
not mentioning $\sigma$), matching the translation clause $\F{a
=_\sigma b} = \compf(\erase{a}) \feq \compf(\erase{b})$; and the
interpretation of a family is index-\emph{aware} --- membership is a
relation between an index tuple and an element --- matching the
$(1{+}m)$-ary guard predicate.

\begin{definition}[Element rank]
\label{def:rank}
For $d$ with $(\tup{d}, d) \in \semr{I}{}$ for some $\tup{d}$, let
$\rank{I}{d}$ be the least $s$ such that $(\tup{d}, d) \in \stage{I}{s}$
for some $\tup{d}$.  (By \cref{lem:index-reflection}(3) the tuple is in
fact unique, but the definition does not need this.)  For $\pv{J}$,
ranks of tuples are defined as in \cite{Extraction2026}: the least $s$
with $\tup{d} \in \pv{J}^s$.
\end{definition}

\begin{lemma}[Index reflection]
\label{lem:index-reflection}
Let $I$ be a datatype as above and $(\tup{d}, d) \in \semr{I}{}$.  Then:
\begin{enumerate}[leftmargin=2em]
\item (\emph{Decomposition.})  There are a constructor index $i$, closed
  $\lbox$ terms $\tup{e} = e_1, \dots, e_{r_i}$ and the valuation
  $\rho_0 = [\tup{y}_i \mapsto \cls{\tup{e}}\,]$ such that $d =
  \cls{\wcns{c}_i(\tup{e})}$; $\cls{e_j} \in \semr{A_{ij}}{\rho_0}$ for
  every non-recursive $j$; $(\sem{\tup{v}_{ij}}_{\rho_0}, \cls{e_j}) \in
  \semr{I}{}$ for every recursive $j$; $\pvr{\varphi_{il}}{\rho_0} = 1$
  for every $l$; and
  \[
    \tup{d} = \sem{\tup{t}_i}_{\rho_0}.
  \]
\item (\emph{Uniqueness.})  The index $i$ and the classes
  $\cls{e_1}, \dots, \cls{e_{r_i}}$ (hence $\rho_0$, hence all data of
  (1)) are uniquely determined by $d$.
\item (\emph{Index uniqueness.})  If also $(\tup{d}', d) \in \semr{I}{}$
  then $\tup{d}' = \tup{d}$.
\item (\emph{Rank descent.})  $\rank{I}{d}$ is defined, and for every
  recursive $j$ in the decomposition of (1),
  $\rank{I}{\cls{e_j}} < \rank{I}{d}$.
\item (\emph{Completeness of stages.})  Conversely, any data as in (1)
  --- an $i$, classes $\cls{\tup{e}}$ with the memberships, premises and
  recursive conditions displayed there --- yield
  $(\sem{\tup{t}_i}_{\rho_0}, \cls{\wcns{c}_i(\tup{e})}) \in \semr{I}{}$.
\end{enumerate}
\end{lemma}

\begin{proof}
(1) $(\tup{d}, d) \in \stage{I}{s+1}$ for some $s$; unfolding
\cref{def:interp} gives exactly the displayed data, with the recursive
conditions in $\stage{I}{s} \subseteq \semr{I}{}$.

(2) Suppose $d = \cls{\wcns{c}_i(\tup{e})} = \cls{\wcns{c}_{i'}
(\tup{e}')}$ arise from two decompositions.  By
\cref{cor:discrimination}(1), $i = i'$ and $\cls{e_j} = \cls{e'_j}$ for
all $j$.  Since all conditions of (1) are conditions on the classes
(interpretations, truth values and denotations are class functions),
the two decompositions coincide.

(3) Apply (1) to $(\tup{d}, d)$ and to $(\tup{d}', d)$; by (2) both
decompositions have the same $i$ and $\rho_0$, so $\tup{d} =
\sem{\tup{t}_i}_{\rho_0} = \tup{d}'$.

(4) Definedness is immediate.  Let $s_0 = \rank{I}{d}$, witnessed by
$(\tup{d}_0, d) \in \stage{I}{s_0}$; necessarily $s_0 = s + 1$ for some
$s$, and unfolding the stage gives a decomposition whose recursive
components satisfy $(\sem{\tup{v}_{ij}}_{\rho_0}, \cls{e_j}) \in
\stage{I}{s}$, whence $\rank{I}{\cls{e_j}} \le s < s_0$.  By (2) this
decomposition is \emph{the} decomposition of (1).

(5) Each recursive condition $(\sem{\tup{v}_{ij}}_{\rho_0}, \cls{e_j})
\in \semr{I}{}$ holds at some finite stage; by monotonicity all hold at
their maximum $s$, and the defining clause puts the conclusion pair in
$\stage{I}{s+1}$.
\end{proof}

Part (1) is the \emph{inversion} direction of fording (used for
hypothesis-position content: the forded inversion axiom, the case rule);
part (5) is the \emph{introduction} direction (constructor typing,
goal-position content); part (3) is the semantic fact that in the term
model every element carries its indices --- \cref{prop:index-unique}
derives its first-order shadow; and parts (1)+(3) together are what
uniqueness of identity proofs buys in type theory, spent here by the
$\tscase{}{}{}{}$ and $\tcast{}{}{}$ cases of \cref{lem:fundamental}
(see \cref{sec:formulas-uip}).

\begin{lemma}[Elementary properties]
\label{lem:sem-basic}
\begin{enumerate}[leftmargin=2em]
\item $\semr{\sigma}{\rho}$, $\pvr{\varphi}{\rho}$, $\kv{P}_\rho$ depend
  only on the restriction of $\rho$ to the free variables of the phrase,
  and are unchanged by extending $\rho$ on fresh variables.
\item Every $(\tup{d}, d) \in \semr{I}{}$ has $(\tup{d},d) \in
  \stage{I}{s}$ for all $s > \rank{I}{d}$; similarly for $\pv{J}$.
\end{enumerate}
\end{lemma}

\begin{proof}
(1) routine induction on the phrase (index terms depend on $\rho$ only
through their free variables).  (2) monotonicity of stages and
\cref{lem:index-reflection}(3): the witnessing tuple at the minimal
stage is the unique tuple.
\end{proof}

\begin{lemma}[Semantic substitution]
\label{lem:sem-subst}
Let $a$ be a term, $\pi$ a proof and $P$ a predicate expression, with
$\rho$ covering the relevant free variables.
\begin{enumerate}[leftmargin=2em]
\item $\sem{\subst{t}{a}{x}}_\rho = \sem{t}_{\rho[x \mapsto
  \sem{a}_\rho]}$, and analogously for substituted types
  ($\semr{\subst{\sigma}{a}{x}}{\rho} = \semr{\sigma}{\rho[x \mapsto
  \sem{a}_\rho]}$), propositions and predicate expressions.
\item Substitution of proofs is invisible:
  $\sem{\subst{t}{\pi}{p}}_\rho = \sem{t}_\rho$,
  $\semr{\subst{\sigma}{\pi}{p}}{\rho} = \semr{\sigma}{\rho}$,
  $\pvr{\subst{\varphi}{\pi}{p}}{\rho} = \pvr{\varphi}{\rho}$.
\item $\pvr{\subst{\varphi}{P}{X}}{\rho} = \pvr{\varphi}{\rho[X \mapsto
  \kv{P}_\rho]}$, and analogously for types and terms (invisible).
\end{enumerate}
\end{lemma}

\begin{proof}
(1) For terms, by \cref{lem:erase-subst} as in \cite{Extraction2026}.
For types, propositions and predicates: induction on the phrase; the
new case $\semr{I(\tup{u})}{\rho}$ follows from the term statement
applied to each index term, since the clause depends on $\tup{u}$ only
through $\sem{\tup{u}}_\rho$.  (2) Erasure discards the substituted
proof (\cref{lem:erase-subst}); proofs do not occur in any semantic
clause --- in particular proofs inside \emph{index terms} are invisible,
which is the model-theoretic form of \cref{rem:proof-indices}.  (3) As
in (1).
\end{proof}

\begin{lemma}[Conversion invariance]
\label{lem:conv-inv}
If $t \conv t'$ then $\sem{t}_\rho = \sem{t'}_\rho$.  If $\sigma \conv
\sigma'$ (both well formed) then $\semr{\sigma}{\rho} =
\semr{\sigma'}{\rho}$; if $\varphi \conv \varphi'$ then
$\pvr{\varphi}{\rho} = \pvr{\varphi'}{\rho}$.
\end{lemma}

\begin{proof}
For terms: $t \conv t'$ implies $\erase{t} \lconv \erase{t'}$
(\cref{lem:erase-sim}), and $\lconv$ is substitutive.  For types and
propositions, consider a single step and proceed by induction on the
structure of the phrase, using \cref{lem:rigid}: the step occurs inside
an embedded term (an index term of a datatype occurrence, an equation
side, a predicate-atom argument) --- covered by the term statement,
since the semantic clauses depend on embedded terms only through
$\sem{\cdot}_\rho$ --- or inside a component type or proposition,
covered by the induction hypothesis at every extension of $\rho$.  The
general claim follows by induction on the conversion path.
\end{proof}

\subsection{The first-order structure \texorpdfstring{$\ML$}{M}}
\label{sec:semantics-model}

\begin{definition}[The structure $\ML$]
\label{def:model}
$\ML$ is the first-order $\SigFOL$-structure with universe $\Dom$ and:
\begin{itemize}[leftmargin=2em]
\item $\feq$ interpreted as identity of classes; $\fapp$ as $\cdot$;
  $\prfc$ as $\cls{\bx}$;
\item $\Ihat{I}^{\ML} = \{ (d, \tup{d}) \mid (\tup{d}, d) \in
  \semr{I}{} \}$ for each datatype $I$ (subject first, matching
  \cref{def:signature}), and $\Ihat{J}^{\ML} = \pv{J}$ for each
  inductive predicate $J$;
\item $\Ihat{\wcns{c}}^{\ML}(d_1, \dots, d_r) = \cls{\wcns{c}(e_1,
  \dots, e_r)}$ for representatives $e_i \in d_i$;
\item for every symbol $f$ introduced by compilation with closure arity
  $s$, extension arity $m$ and witness $e_f$ (free variables $\tup{w}$):
  writing $\tup{d}$ for the values at the closure places and $\tup{d}'$
  for those at the extension places,
  $f^{\ML}(\dots) = \cls{\subst{e_f}{\tup{e}}{\tup{w}}\;
  e'_1 \cdots e'_m}$ for representatives.  The argument places appear
  in the order of the symbol's introduction in \cref{def:compC}:
  closure places first for constants and for $\lambda$- and fix-lifts,
  whereas a case-lift symbol $g(z, \tup{w}')$ takes its single
  extension place (the scrutinee) first, so for it
  $g^{\ML}(d, \tup{d}') =
  \cls{\subst{e_g}{\tup{e}'}{\tup{w}'}\;e}$.  In particular
  $\Ihat{c}^{\ML}(\tup{d}) = \cls{c\,\tup{e}}$ and $(\Ihat{c}^0)^{\ML} =
  \cls{c}$.
\end{itemize}
All interpretations are well defined on classes by congruence of
$\lconv$.  We write $\den{u}{\rho}$ for the denotation of a first-order
term $u$ under an assignment $\rho$, and $\ML, \rho \vDash \phi$ for
Tarskian satisfaction.  Valuations restricted to term variables serve as
first-order assignments.
\end{definition}

\subsection{Adequacy of the term translation and of compilation}
\label{sec:semantics-adequacy}

The two lemmas of this subsection are the companion paper's, re-proved
here in full --- including the cases for the two head clauses added in
\cref{def:compC} --- both because the soundness proof cites their
content and because their \emph{independence} of the indexed extension
is itself the point: compilation concerns only the program half, which
is index-blind.

\begin{definition}[Coherent template maps]
\label{def:coherent}
For a template map $\theta$ and a substitution $\sigma_\theta$ assigning
closed $\lbox$ terms to $\dom{\theta}$, say $(\theta, \sigma_\theta)$ is
\emph{coherent} for an assignment $\rho$ if for every $g$ with
$\theta(g) = (f, \tup{w}, m)$ and all $d_1, \dots, d_m \in \Dom$ with
representatives $e_i$:
$f^{\ML}(\den{\tup{w}}{\rho}, \tup{d}) =
\cls{\sigma_\theta(g)\;e_1 \cdots e_m}$.
\end{definition}

\begin{lemma}[Adequacy of $\compf$]
\label{lem:comp-adequacy}
Let $e$ be a $\lbox$ term whose free variables are first-order variables
assigned by $\rho$ or bound by $\theta$, and let $(\theta,
\sigma_\theta)$ be coherent for $\rho$.  Then $\den{\compf_\theta(e)}
{\rho} = \cls{e^{\rho, \sigma_\theta}}$, where $e^{\rho,\sigma_\theta}$
substitutes representatives of $\rho(x)$ for the first-order variables
and $\sigma_\theta(g)$ for the $\theta$-bound variables.  In particular
$\den{\compf(\erase{t})}{\rho} = \sem{t}_\rho$ for every $\CICi$ term
$t$ with $\theta = \emptyset$.
\end{lemma}

\begin{proof}
By induction on $e$ following the clauses of \cref{def:compC}.  Fix
representatives $\hat{e}_i \in \den{\compf_\theta(e_i)}{\rho}$; by the
induction hypothesis we may take $\hat{e}_i = e_i^{\rho,
\sigma_\theta}$.  Since $\fapp^{\ML}$ is the application of classes,
any $\fapp$-chain below denotes the class of the corresponding applied
representatives.
\begin{itemize}[leftmargin=2em]
\item $e = x$ ordinary: $\den{x}{\rho} = \rho(x) =
  \cls{x^{\rho,\sigma_\theta}}$ by the definition of the substitution.
  $e = \bx$: both sides are $\cls{\bx}$.  $e = \wcns{c}(\tup{e})$:
  $\den{\Ihat{\wcns{c}}(\compf_\theta(\tup{e}))}{\rho} =
  \Ihat{\wcns{c}}^{\ML}(\cls{\tup{e}^{\rho,\sigma_\theta}}) =
  \cls{\wcns{c}(\tup{e}^{\rho,\sigma_\theta})}$ by the induction
  hypothesis and \cref{def:model}.
\item Spine with $\theta$-bound head, $\theta(g) = (f, \tup{w}, m)$:
  the denotation of $f(\tup{w}, \compf_\theta(e_1), \dots,
  \compf_\theta(e_m))$ is $f^{\ML}(\den{\tup{w}}{\rho},
  \cls{\hat{e}_1}, \dots, \cls{\hat{e}_m}) =
  \cls{\sigma_\theta(g)\,\hat{e}_1 \cdots \hat{e}_m}$ by coherence
  (\cref{def:coherent}); the $\fapp$-chain appends $\hat{e}_{m+1},
  \dots, \hat{e}_j$, giving $\cls{\sigma_\theta(g)\,\hat{e}_1 \cdots
  \hat{e}_j} = \cls{e^{\rho,\sigma_\theta}}$.
\item Constant head $c$ with $j \geq a_c$:
  $\Ihat{c}^{\ML}(\cls{\hat{e}_1}, \dots, \cls{\hat{e}_{a_c}}) =
  \cls{c\,\hat{e}_1 \cdots \hat{e}_{a_c}}$ by \cref{def:model} (the
  witness of $\Ihat{c}$ is the constant $c$), and the $\fapp$-chain
  appends the rest.  With $j < a_c$: $(\Ihat{c}^0)^{\ML} = \cls{c}$ and
  the $\fapp$-chain gives $\cls{c\,\hat{e}_1 \cdots \hat{e}_j}$.
\item Variable head: $\rho(x) \cdot \cls{\hat{e}_1} \cdots
  \cls{\hat{e}_j} = \cls{x^{\rho}\,\hat{e}_1 \cdots \hat{e}_j}$.
\item $\bx$-head: $\prfc^{\ML} = \cls{\bx}$ and the $\fapp$-chain gives
  $\cls{\bx\,\hat{e}_1 \cdots \hat{e}_j}$, the class of
  $e^{\rho,\sigma_\theta}$ --- dead code applied is junk, but a
  perfectly good domain element.
\item Constructor head $\wcns{c}(\tup{e}')$: by the constructor case,
  $\den{\Ihat{\wcns{c}}(\compf_\theta(\tup{e}'))}{\rho} =
  \cls{\wcns{c}(\tup{e}'^{\rho,\sigma_\theta})}$, and the
  $\fapp$-chain appends $\hat{e}_1, \dots, \hat{e}_j$.
\item Complex heads, via their value translations.  $\lambda$-lift $h$
  with value $f(\tup{w})$: by \cref{def:model}, $f^{\ML}
  (\den{\tup{w}}{\rho}) = \cls{\subst{h}{\tup{\hat{w}}}{\tup{w}}} =
  \cls{h^{\rho,\sigma_\theta}}$, the witness being $h$ itself and the
  closure parameters receiving the values $\rho$ assigns them.
  Case-lift
  $h = \mathsf{case}(s; \tup{br})$ with value $g(\compf_\theta(s),
  \tup{w}')$: $g^{\ML}(\cls{\hat{s}}, \den{\tup{w}'}{\rho}) =
  \cls{(\lambda z.\mathsf{case}(z; \tup{br}))
  [\tup{\hat{w}}'/\tup{w}']\;\hat{s}}$, which $\beta$-reduces in one
  step to $\cls{\mathsf{case}(\hat{s};
  \tup{br}[\tup{\hat{w}}'/\tup{w}'])} = \cls{h^{\rho,\sigma_\theta}}$,
  using the induction hypothesis for $\hat{s}$.  Fix-lift with value
  $f^0(\tup{w})$: $(f^0)^{\ML}(\den{\tup{w}}{\rho}) =
  \cls{\subst{h}{\tup{\hat{w}}}{\tup{w}}} = \cls{h^{\rho,
  \sigma_\theta}}$ via the witness $e_{f^0} = h$.  In each case the
  $\fapp$-chain then appends $\hat{e}_1, \dots, \hat{e}_j$.
\end{itemize}
The final claim is the special case $e = \erase{t}$ with $\theta =
\emptyset$: free variables of $\erase{t}$ are free term variables of
$t$, so $\cls{\erase{t}^{\rho,\emptyset}} = \cls{\erase{t}^{\rho}} =
\sem{t}_\rho$ by \cref{def:valuation}.
\end{proof}

\begin{lemma}[Soundness of compilation]
\label{lem:compile-sound}
Consider any invocation of the compilation procedure on the environment
or on a term embedded in the goal or in a declaration, and let
$\mathrm{B}_\theta(L; e)$ range over the body-compilation calls it
performs.  Then:
\begin{enumerate}[leftmargin=2em]
\item for every such call there is a substitution $\sigma_\theta$ such
  that for every assignment $\rho$ of the variables of $L$, the pair
  $(\theta, \sigma_\theta)$ is coherent for $\rho$ and $\den{L}{\rho} =
  \cls{e^{\rho, \sigma_\theta}}$;
\item every axiom emitted by the procedure holds in $\ML$ --- under
  \emph{every} assignment of its universally quantified variables, junk
  included.
\end{enumerate}
\end{lemma}

\begin{proof}
By induction on the tree of procedure calls; part (1) is the invariant
that makes part (2) a local computation at each emitting step.

\emph{Roots.}  For $c \defeq t : \tau$ with $\erase{t} = \lambda
\tup{z}.e_0$ (maximal prefix): $\theta = \emptyset$, and for any
$\rho$ with representatives $\tup{e}$ of $\rho(\tup{z})$,
$\den{\Ihat{c}(\tup{z})}{\rho} = \cls{c\,\tup{e}}$
(\cref{def:model}) $= \cls{\erase{t}\,\tup{e}}$ (one $\delta$-step)
$= \cls{e_0^\rho}$ ($|\tup{z}|$ $\beta$-steps).  For $\erase{t} =
\mathsf{fix}\,g_0.\lambda\tup{z}.e_0$: set $\sigma_\theta(g_0) \defeq
\erase{t}$; coherence of $(\theta, \sigma_\theta)$ is the computation
$\Ihat{c}^{\ML}(\tup{d}) = \cls{c\,\tup{e}} =
\cls{\erase{t}\,\tup{e}} = \cls{\sigma_\theta(g_0)\,\tup{e}}$, and the
invariant for the initial call $\mathrm{B}(\Ihat{c}(\tup{z}); e_0)$
follows by one further $\mu$-step and the $\beta$-steps.  For a
$\lambda$-lift ($\mathrm{B}_\theta(f(\tup{w}); h)$) the invariant is
the interpretation of $f$ through its witness; for a case-lift
($\mathrm{B}_\theta(g(z, \tup{w}'); \mathsf{case}(z;\tup{br}))$) it is
the witness interpretation followed by one $\beta$-step ($e_g$ is the
$\lambda$-wrapped case); for a fix-lift, extend $\sigma_\theta$ by
mapping $g_0$ to the \emph{closed instance} of the witness at
representatives of $\den{\tup{w}}{\rho}$ --- coherence for the new
binding is then the same computation as at the root, one level down.

\emph{Rule (1)} ($L \mapsto L \fapp x$, $e = \lambda x.e'$): an
assignment of the extended left-hand side restricts to one of $L$; the
invariant transfers by one $\beta$-step on representatives,
$\cls{(\lambda x.e'^{\rho,\sigma_\theta})\,\hat{x}} =
\cls{e'^{\rho[x],\sigma_\theta}}$.

\emph{Rule (2)} (refinement of a pattern variable): an assignment
$\rho'$ of the variables of
$\subst{L}{\Ihat{\wcns{c}}_i(\tup{z}_i)}{x}$ induces the assignment
$\rho$ of $L$ that agrees with $\rho'$ elsewhere and maps $x$ to
$\Ihat{\wcns{c}}_i^{\ML}(\rho'(\tup{z}_i)) =
\cls{\wcns{c}_i(\tup{e})}$; the two left-hand sides then denote the
same element, and on the right, a head $\iota$-step converts the
representative of $\mathsf{case}(x;\dots)^{\rho,\sigma_\theta}$ into
that of $(\subst{b_i}{\wcns{c}_i(\tup{z}_i)}{x})^{\rho',
\sigma_\theta}$.  Note that the $\iota$-step of $\lbox$ fires on
\emph{any} constructor form, with no index side condition --- this is
why the emitted equations hold at all elements, wrong-index
instantiations included, which \cref{prop:match-equations} isolates.

\emph{Rule (3)}: by the invariant, the left-hand side of the emitted
axiom denotes $\cls{\mathsf{case}(s;\tup{br})^{\rho,\sigma_\theta}}$;
by \cref{lem:comp-adequacy} and one $\beta$-step through the witness
of $g$, so does the right-hand side $g(\compf_\theta(s), \tup{w}')$;
and the invariant for the recursive call on $g(z, \tup{w}')$ is the
case-lift root computation above.  \emph{Rule (4)}: both sides denote
$\cls{e^{\rho,\sigma_\theta}}$, the left by the invariant and the
right through the fix-lift witness of $f^0$.  \emph{Rule (5)}: the
invariant identifies the left-hand side's denotation with
$\cls{e^{\rho,\sigma_\theta}}$, and \cref{lem:comp-adequacy}
identifies the right-hand side's with the same class.  \emph{Pointer
axioms}: for a constant, both sides of $\Ihat{c}^0 \fapp z_1 \fapp
\cdots \fapp z_{a_c} \feq \Ihat{c}(\tup{z})$ denote
$\cls{c\,\tup{e}}$; for a fix-lift, both sides denote the witness
instance applied to the extension arguments.
\end{proof}

\subsection{Coincidence: the translation reflects the semantics}
\label{sec:semantics-coincidence}

\begin{lemma}[Coincidence]
\label{lem:coincidence}
Let $\sigma$ be a shallow set type and $\varphi$ a shallow proposition
over $\Env$ and a context of term variables, and let $\rho$ be a
valuation of their free variables.  Then:
\begin{enumerate}[leftmargin=2em]
\item for every first-order term $u$ over the variables of $\rho$:\quad
  $\ML, \rho \vDash \guard{\sigma}{u} \iff \den{u}{\rho} \in
  \semr{\sigma}{\rho}$;
\item $\ML, \rho \vDash \F{\varphi} \iff \pvr{\varphi}{\rho} = 1$.
\end{enumerate}
\end{lemma}

\begin{proof}
Mutual induction on the structure of $\sigma$ and $\varphi$, following
\cref{def:guard-F}.
\begin{itemize}[leftmargin=2em]
\item $\sigma = I(\tup{a})$: $\ML, \rho \vDash \Ihat{I}(u,
  \compf(\erase{\tup{a}}))$ iff $(\den{u}{\rho},
  \den{\compf(\erase{\tup{a}})}{\rho}) \in \Ihat{I}^{\ML}$ iff
  $(\sem{\tup{a}}_\rho, \den{u}{\rho}) \in \semr{I}{}$ --- using
  \cref{lem:comp-adequacy} to identify
  $\den{\compf(\erase{a_h})}{\rho} = \sem{a_h}_\rho$ --- iff
  $\den{u}{\rho} \in \semr{I(\tup{a})}{\rho}$ by \cref{def:interp}.
\item $\sigma = \Pi x{:}\sigma'.\tau$ (with $x$ $\alpha$-renamed fresh
  for $u$ and $\rho$): $\ML, \rho \vDash \forall x.\ \guard{\sigma'}{x}
  \to \guard{\tau}{u \fapp x}$ iff for every $d \in \Dom$, $\ML,
  \rho[x \mapsto d] \vDash \guard{\sigma'}{x}$ implies $\ML, \rho[x
  \mapsto d] \vDash \guard{\tau}{u \fapp x}$; by both induction
  hypotheses (at the extended valuation) this says: for every $d \in
  \semr{\sigma'}{\rho}$, $\den{u}{\rho} \cdot d \in
  \semr{\tau}{\rho[x \mapsto d]}$ --- which is the defining condition
  of $\den{u}{\rho} \in \semr{\Pi x{:}\sigma'.\tau}{\rho}$
  (\cref{def:interp}).
\item $\sigma = \varphi' \to \tau$: satisfaction of $\F{\varphi'} \to
  \guard{\tau}{u}$ says $\pvr{\varphi'}{\rho} = 1$ implies
  $\den{u}{\rho} \in \semr{\tau}{\rho}$ (both induction hypotheses),
  the defining condition of $\semr{\varphi' \to \tau}{\rho}$.
\item $\sigma = \refty{x{:}\sigma'}{\varphi'}$: the first conjunct is
  the induction hypothesis for $\sigma'$; for the second, the
  first-order substitution lemma gives $\ML, \rho \vDash
  \subst{\F{\varphi'}}{u}{x}$ iff $\ML, \rho[x \mapsto \den{u}{\rho}]
  \vDash \F{\varphi'}$, which by induction hypothesis (2) is
  $\pvr{\varphi'}{\rho[x \mapsto \den{u}{\rho}]} = 1$; together, the
  defining condition of the refinement interpretation.
\item Connectives, $\top$, $\bot$: $\Ff$ is homomorphic and both
  $\vDash$ and $\pv{\cdot}$ interpret the connectives by the classical
  truth tables, so induction hypothesis (2) on the components
  concludes.
\item Quantifiers: $\ML, \rho \vDash \forall x.\ \guard{\sigma'}{x}
  \to \F{\varphi'}$ iff $\pvr{\varphi'}{\rho[x \mapsto d]} = 1$ for
  every $d \in \semr{\sigma'}{\rho}$ (both induction hypotheses, as in
  the $\Pi$-case) iff $\min_{d \in \semr{\sigma'}{\rho}}
  \pvr{\varphi'}{\rho[x \mapsto d]} = 1 = \pvr{\forall
  x{:}\sigma'.\varphi'}{\rho}$; the empty-range conventions agree
  (vacuous satisfaction, $\min\emptyset = 1$).  Dually for $\exists$
  with $\wedge$ and $\max\emptyset = 0$.
\item Equations and $J$-atoms: by \cref{lem:comp-adequacy},
  $\den{\compf(\erase{a})}{\rho} = \sem{a}_\rho$, and the model
  interprets $\feq$ as class equality and $\Ihat{J}$ as $\pv{J}$.
\end{itemize}
Shallowness is used exactly here: predicate quantifiers and
predicate-variable atoms have no translation; all remaining formers are
covered.
\end{proof}

The new clause deserves a comment in the direction of
\cref{sec:fording-obligations}: the coincidence for family guards is an
equivalence \emph{at every index tuple}, junk tuples included --- both
sides are false at a wrong-index pair such as $(\cls{\csvnil},
\cls{\csS(\csO)})$ by \cref{lem:index-reflection}(3).  It is this
two-sided, index-aware coincidence that licenses using guard atoms in
\emph{both} polarities of the translation, and its failure under
index-forgetting that \cref{prop:index-dropping} exhibits.
